% SOURCE_AUTHORITY: sga5_fr_workpass.tex, lines 11268--11289 (LF-normalized unit).
% SOURCE_SHA256: 791F4EFFC5E02832D5D77ED03518C8156D6F07E4C8238B03545DB93D883FBB28
Passemos agora à «fórmula-chave» mencionada na introdução. As notações adotadas, em grande parte já recordadas, são as de 4. Observe-se que são compatíveis com as introduzidas no presente parágrafo, salvo que o esquema denotado por $Y'$ em 4 é denotado aqui por $H$ (por «hiperplano»). Utilizaremos, de resto, ambas as notações indistintamente.

\begin{statement}{Proposição 9.6}
Sejam $k$ um corpo separadamente fechado e $i:Y\hookrightarrow X$ uma imersão fechada, de codimensão pura $d$, de $k$-esquemas quase projetivos e lisos. Denotando por $f:X'\to X$ o $X$-esquema obtido pela explosão de $X$ ao longo de $Y$, seja
\[
\begin{tikzcd}
Y'\arrow[r,"j"]\arrow[d,"g"']&X'\arrow[d,"f"]\\
Y\arrow[r,"i"']&X
\end{tikzcd}
\]
o diagrama cartesiano construído a partir de $f$ e $i$. Então, denotando por $F$ o $\mathcal O_{Y'}$-módulo localmente livre definido pela sequência exata canônica
\[
0\to F\to g^*(N)\to\mathcal O_{Y'}(1)\to0,
\tag{9.6.1}
\]
tem-se a igualdade
\[
f^*i_*(y)=j_*(g^*(y)c_{d-1}(\check F))
\tag{9.6.2}
\]
em $C^*(X')$, para todo $y\in C^*(Y)$.
\end{statement}
